\title{QLoRA: Efficient Finetuning of Quantized LLMs}

\begin{abstract}
We introduce \textsc{QLoRA}, a novel method for efficient finetuning that substantially decreases memory consumption, enabling the finetuning of a 65B parameter language model on a single 48GB GPU while matching the task performance of traditional 16-bit finetuning. With \textsc{QLoRA}, gradients are propagated through a frozen 4-bit quantized pretrained language model and updated within Low Rank Adapters~(LoRA). Our leading set of models, named \textbf{Guanaco}, surpasses all prior openly accessible models on the Vicuna benchmark, achieving 99.3\% of ChatGPT’s performance after just 24 hours of finetuning on a single GPU. \textsc{QLoRA} incorporates several technical contributions to minimize memory requirements without loss of efficacy: (a) 4-bit NormalFloat (NF4), a new data format that is theoretically optimal for representing weights drawn from a normal distribution, (b) Double Quantization, which further compresses model memory usage by quantizing the quantization parameters themselves, and (c) Paged Optimizers, which regulate and mitigate memory usage fluctuations. We apply \textsc{QLoRA} to finetune over 1,000 distinct models, thoroughly assessing instruction-following and conversational capabilities across 8 instruction datasets, a variety of architectures (LLaMA, T5), and scales beyond what is practical with conventional methods (e.g., 33B and 65B parameter models). Our experiments reveal that using QLoRA finetuning on compact, high-quality datasets attains state-of-the-art results, sometimes even outperforming earlier benchmarks with smaller models. Comprehensive evaluation, including both human assessment and GPT-4-based metrics, suggests that GPT-4 evaluations serve as a low-cost, reliable proxy for human judgment. Moreover, we observe significant shortcomings in the reliability of current chatbot evaluation benchmarks. Our ‘lemon-picked’ analysis highlights failure modes where \textbf{Guanaco} does not match ChatGPT. We are releasing all our models and supporting code, with CUDA kernels facilitating 4-bit training.\footnote{\url{https://github.com/artidoro/qlora} and \url{https://github.com/TimDettmers/bitsandbytes}}
\end{abstract}

\section{Introduction}

Finetuning large language models (LLMs) has proven to be a powerful strategy for enhancing model performance \citep{min2021metaicl, wei2021finetuned, ouyang2022training, wang2022super, wang2022self, liu2022few}, as well as for incorporating preferred behaviors or eliminating unwanted ones \citep{ouyang2022training,askell2021general,bai2022training}. Despite its benefits, finetuning very large models remains prohibitively costly; for instance, standard 16-bit finetuning of a LLaMA model with 65B parameters~\citep{touvron2023llama} necessitates more than 780 GB of GPU memory. Although contemporary quantization techniques can help minimize LLMs' memory consumption \citep{dettmers2022llm, dettmers2022case, frantar2022gptq, xiao2022smoothquant}, these methods are primarily designed for inference and are unsuitable for use during training, as their effectiveness deteriorates \citep{wortsman2023stable}.

In this work, we show for the first time that finetuning a quantized model at 4-bit precision can be accomplished with no loss in performance. Our proposed approach, \textsc{QLoRA}, employs an innovative high-precision method to quantize the pretrained model to 4-bit and then integrates a compact set of trainable Low-rank Adapter weights \citep{hu2021lora}, 

which are optimized by propagating gradients through the quantized parameters.

With \textsc{QLoRA}, the average GPU memory required to finetune a 65B parameter model drops from over 780GB to under 48GB, all while maintaining both runtime efficiency and predictive accuracy on par with conventional 16-bit full finetuning. This dramatically lowers the resource barrier for LLM finetuning, enabling the largest open models to be finetuned using a single GPU. Utilizing \textsc{QLoRA}, we introduce the \textbf{Guanaco} series of models: our second-best model achieves 97.8\% of ChatGPT's performance on the Vicuna~\citep{vicuna2023} evaluation, training in under 12 hours on a standard consumer GPU. By deploying a single professional GPU for 24 hours, we reach 99.3\% with our largest model, effectively matching ChatGPT on the Vicuna benchmark. In practical deployment, our smallest \textbf{Guanaco} model (7B parameters) operates with just 5 GB of memory and surpasses the 26 GB Alpaca model by more than 20 percentage points on the Vicuna benchmark (see Table~\ref{tbl:vicuna_eval}).

\textsc{QLoRA} presents several novel strategies aimed at minimizing memory requirements while maintaining high levels of model performance: (1) \textbf{4-bit NormalFloat}, a quantization scheme theoretically optimal for normally distributed variables, demonstrates superior empirical results compared to traditional 4-bit Integers and Floats. (2) \textbf{Double Quantization}, an approach that quantizes the quantization coefficients themselves, reducing average memory consumption by about 0.37 bits per parameter (equivalent to roughly 3 GB for a 65B parameter model). (3) \textbf{Paged Optimizers}, which employ NVIDIA unified memory to circumvent the abrupt increases in memory demand caused by gradient checkpointing when handling mini-batches with extended sequence lengths. These techniques are integrated within an enhanced LoRA framework, incorporating adapters at every layer of the network, which effectively mitigates the accuracy compromises observed in previous studies.

Owing to the increased memory efficiency of \textsc{QLoRA}, it becomes feasible to conduct comprehensive evaluations of instruction finetuning methods and chatbot performance across model scales that are prohibitive for standard finetuning procedures because of their memory footprints. To this end, we have trained over 1,000 models spanning a range of instruction tuning datasets, model architectures, and parameter sizes from 80M up to 65B. Beyond demonstrating that \textsc{QLoRA} achieves parity with 16-bit training performance (\S\ref{sec:comp_to_std}) and enabling the training of a leading-edge chatbot, \textbf{Guanaco}, (\S\ref{sec:pushing}), we further investigate emergent patterns in model performance. Our analysis reveals, first, that dataset quality substantially outweighs sheer size—for example, a 9,000-example dataset (OASST1) outperformed a subsampled 450,000-example dataset (FLAN v2) in chatbot benchmarks, even when both were crafted for instruction-following tasks. Second, we show that high scores on the Massive Multitask Language Understanding (MMLU) benchmark do not necessarily translate to superior performance on the Vicuna chatbot benchmark, and vice versa, highlighting the centrality of task-specific dataset appropriateness over mere scale.

In addition, we conduct a comprehensive evaluation of chatbot performance by leveraging assessments from both human annotators and GPT-4. Our approach employs a tournament-style benchmark, wherein different models are pitted against one another to generate optimal responses for the same prompt. Each contest is adjudicated either by GPT-4 or human raters, who select the superior answer. The cumulative tournament outcomes are summarized via Elo scores~\citep{elo1967proposed, elo1978rating}, which serve to rank the competing chatbots. Our findings indicate substantial concordance between human judgments and GPT-4 rankings regarding overall model performance, although notable discrepancies do occur in certain cases. These observations underscore that model-based evaluations, while cost-effective relative to human annotation, can still introduce significant uncertainties.

To further contextualize the quantitative results, we supplement our chatbot benchmarking with an in-depth qualitative examination of the \textbf{Guanaco} models. Through this analysis, we identify representative instances of both successes and failures that elude capture in standardized quantitative assessments.

All generated responses, alongside both human and GPT-4 ratings, are publicly released to support future research. We have also open-sourced our complete codebase, associated CUDA kernels, and provided seamless integration with the Hugging Face transformers library~\citep{wolf2019huggingface}. Furthermore, we distribute adapter modules for models of sizes 7/13/33/65B, each trained across eight distinct instruction-tuning datasets—yielding a suite of 32 different fine-tuned, open-source models.

\section{Background}
\paragraph{Block-wise k-bit Quantization} Quantization refers to the technique of mapping a data representation with higher information content to one with reduced information, commonly achieved by converting higher bit-width types (such as 32-bit floats) to those of lower bit-width (like 8-bit integers). To maximize the effective utilization of the entire dynamic range in the target low-bit data type, inputs are generally rescaled by normalizing to the absolute maximum value among the input’s elements—typically structured as a tensor. For instance, converting a 32-bit floating-point tensor (FP32) into an 8-bit integer (Int8) tensor with permissible values in $[-127, 127]$ involves:

\begin{equation}
    \mathbf{X}^{ \text{Int8}} = \text{round}\left( \frac{127}{{\text{absmax}}(\mathbf{X}^{{\text{FP32}}})}\mathbf{X}^{{\text{FP32}}} \right) = \text{round}(c^{\text{FP32}}\cdot \mathbf{X}^{{\text{FP32}}}),
\end{equation}

where $c$ denotes the {\em quantization constant} or {\em quantization scale}. The original values can be recovered through dequantization, defined as:

\begin{equation}
    \text{dequant}(c^{\text{FP32}}, \mathbf{X}^{\text{Int8}}) = \frac{\mathbf{X}^{\text{Int8}}}{c^{\text{FP32}}} = \mathbf{X}^{\text{FP32}}
\end{equation}

A drawback of this method arises when an input tensor contains a value with large magnitude, i.e., an outlier; under such circumstances, the quantization bins—each representing distinct bit patterns—are inefficiently filled, with some bins containing few or no values after quantization. To address the problem posed by outliers, a widely used solution is to divide the input tensor into several blocks, each of which is quantized independently and assigned its own quantization constant $c$. Specifically, the input tensor $\mathbf{X}\in \mathbb{R}^{b\times h}$ is segmented into $n$ consecutive blocks of length $B$. This is done by flattening $\mathbf{X}$ and partitioning the resulting sequence into $n = ({b\times h})/{B}$ segments. Each segment undergoes quantization separately, using Equation~1, resulting in a quantized output as well as a distinct quantization constant $c_i$ per block.

\paragraph{Low-rank Adapters} Low-rank Adapter (LoRA) finetuning \citep{hu2021lora} achieves reductions in memory usage by introducing a compact set of trainable parameters (called adapters) while keeping the original model weights frozen during optimization. In this setup, gradients computed during stochastic gradient descent propagate through the fixed backbone weights to update only the adapter parameters, thereby minimizing the loss. LoRA modifies a linear mapping by incorporating an additional factorized projection. If we start with a linear projection ${\mathbf{X} \mathbf{W} = \mathbf{Y}}$, where $\mathbf{X}\in  \mathbb{R}^{b\times h}$ and $\mathbf{W}\in  \mathbb{R}^{h\times o}$, the LoRA transformation is defined as:

\begin{equation}
    \mathbf{Y} = \mathbf{X}\mathbf{W} + s\mathbf{X}\mathbf{L}_1\mathbf{L}_2,
\end{equation}

with $\mathbf{L}_1\in  \mathbb{R}^{h\times r}$, $\mathbf{L}_2\in  \mathbb{R}^{r\times o}$, and a scalar $s$.

\paragraph{Memory Requirement of Parameter-Efficient Finetuning} The issue of memory consumption when training with LoRA, specifically regarding both the quantity and size of adapters, warrants careful consideration. Given LoRA's exceptionally small memory overhead, it is possible to deploy additional adapters to boost performance, incurring minimal increases in total memory usage. Despite LoRA's original purpose as a Parameter Efficient Finetuning (PEFT) strategy, it is important to note that the primary contributor to memory usage during LLM finetuning is the storage of activation gradients, rather than the parameters introduced by LoRA. For instance, in the context of finetuning a 7B LLaMA model on FLAN v2 data with a batch size of 1—where LoRA weights constitute the commonly adopted 0.2\% of the original parameters~\citep{hu2021lora,liu2022few}—the memory required for LoRA input gradients amounts to 567 MB, while the parameters themselves occupy a mere 26 MB. Applying gradient checkpointing~\citep{chen2016training} further reduces average input gradient memory to just 18 MB per sequence, yet this remains greater than the total memory of all LoRA parameters. For context, a 4-bit quantized base model occupies 5,048 MB of memory. This underscores both the role of gradient checkpointing in curbing memory usage and the limited gains from further shrinking the size of LoRA parameters. Thus, one can increase the number of adapters without substantially elevating the overall training memory demand (refer to Appendix~\ref{app:memory} for comprehensive details). As will be elaborated later, this flexibility is key to achieving the performance levels characteristic of full 16-bit precision.

\section{\textsc{QLoRA} Finetuning}

\textsc{QLoRA} enables accurate 4-bit finetuning through two novel approaches we introduce: 4-bit NormalFloat (NF4) quantization and Double Quantization. In addition, we present Paged Optimizers as a solution to avoid the memory spikes during gradient checkpointing, which have historically led to out-of-memory issues and thus hindered large-scale model finetuning on a single machine.

Within \textsc{QLoRA}, a single low-precision storage type is used—typically 4-bit—while computations are carried out using another type, most often BFloat16. Operationally, every time a \textsc{QLoRA} weight tensor is utilized, it is first dequantized to BFloat16 before any 16-bit matrix multiplication is executed.

Below, we elaborate on each component of \textsc{QLoRA} and subsequently offer a formal specification of the method.

\paragraph{4-bit NormalFloat Quantization} The NormalFloat (NF) format extends Quantile Quantization \citep{dettmers2022optimizers}, an information-theoretic optimal approach that distributes input tensor values uniformly among quantization bins. Quantile quantization assigns equal numbers of values from the input tensor to each bin, using the empirical cumulative distribution function to estimate quantiles.

However, the quantile estimation process in quantile quantization is computationally demanding. Therefore, efficient quantile approximation methods, such as SRAM quantiles~\citep{dettmers2022optimizers}, are typically employed. These approximation strategies, due to their inherent limitations, may result in large quantization errors for outlier values—elements that often play a significant role.

Nevertheless, when the input tensors are drawn from a distribution fixed up to a quantization constant, both costly quantile estimation and large approximation errors can be circumvented. In such circumstances, identical quantiles among input tensors allow for precise quantile estimation to be tractable.

Given that pretrained neural network weights typically follow a zero-mean normal distribution with standard deviation $\sigma$ (refer to Appendix~\ref{app:normality}), it is possible to map all such weights onto a uniform reference distribution by appropriately adjusting $\sigma$ so the distribution aligns perfectly with the boundaries of our chosen data type. We define the data type’s range arbitrarily as $[-1, 1]$. Therefore, it is necessary to normalize both the data type quantiles and the neural network weights within this fixed interval.

For zero-mean normal distributions of any standard deviation $\sigma$ within $[-1, 1]$, the information-theoretic optimal data type is determined as follows: (1) derive the $2^k+1$ quantiles of a standard normal distribution $N(0, 1)$ to construct a quantile-based data type with $k$-bits tailored for normal distributions; (2) transform the quantile values of this data type to fit the $[-1, 1]$ range; (3) apply quantization to the input weight tensor by rescaling its values into the $[-1, 1]$ interval using absolute maximum normalization.

When the span of the weights and the data type coincide, quantization can proceed as per usual practice. Notably, the normalization step in (3) corresponds to scaling the standard deviation of the weights to align with the standard deviation of the $k$-bit data type. Formally, the $2^k$ data type values $q_i$ are estimated as:

\begin{equation}
q_i  = \frac{1}{2}\left( Q_X\left(\frac{i}{2^k+1}\right) + Q_X\left(\frac{i+1}{2^k+1}\right)\right),
\end{equation}

where $Q_X(\cdot)$ denotes the quantile function corresponding to the standard normal distribution $N(0,1)$. One notable limitation of symmetric $k$-bit quantization is the absence of a precise zero representation, which is critical for exactly encoding padding and other elements whose value is zero. To guarantee a discrete zero-point at $0$ while also fully utilizing all $2^k$ bit patterns available in a $k$-bit data type, we propose an asymmetric approach: we estimate the quantiles $q_i$ separately for the negative and positive domains—using $2^{k-1}$ quantiles for the negative portion, and $2^{k-1} + 1$ quantiles for the positive side. These two sets of $q_i$ are then merged, and one of the two overlapping zero points is removed to avoid duplication. The resulting data type is referred to as \emph{k-bit NormalFloat} (NFk), characterized by having each quantization interval contain, on expectation, the same number of values; this construction is information-theoretically optimal for normally distributed data centered at zero. A comprehensive list of the explicit values represented by this data type can be found in Appendix~\ref{app:nf4}.

\paragraph{Double Quantization{}} 
We propose \emph{Double Quantization} (DQ), a scheme in which the quantization constants themselves are quantized, providing further reductions in memory usage. Precise 4-bit quantization schemes generally rely on a small blocksize~\citep{dettmers2022case}, which, however, introduces significant overhead for storing quantization constants. As an illustration, with 32-bit quantization constants and a blocksize of 64 for $\mathbf{W}$, this overhead amounts to $32/64 = 0.5$ bits per parameter, on average. Double Quantization addresses this by compressing the quantization constants themselves.

More concretely, Double Quantization applies a second quantization to the first-stage quantization constants $c_2^{\text{FP32}}$, treating them as input data for this subsequent quantization. This produces a new set of quantized quantization constants $c_2^{\text{FP8}}$ along with new quantization constants $c_1^{\text{FP32}}$ at the second quantization level. For this second quantization, we employ 8-bit floating-point types with a blocksize of 256, since empirical evidence~\citep{dettmers2022case} indicates there is no loss in performance at 8-bit precision. Given that $c_2^{\text{FP32}}$ values are strictly positive, we subtract their mean prior to quantization in order to center their distribution and take advantage of symmetric quantization. This strategy, with a blocksize of 64, decreases the storage required per parameter from $32/64 = 0.5$ bits to $8/64 + 32/(64 \cdot 256) = 0.127$ bits, yielding a per-parameter reduction of 0.373 bits.

\paragraph{Paged Optimizers} leverage the NVIDIA unified memory~\footnote{\tiny \url{https://docs.nvidia.com/cuda/cuda-c-programming-guide}} capability, which transparently manages the transfer of memory pages between the CPU and GPU. This automatic paging allows computations on the GPU to continue seamlessly even in circumstances where GPU memory becomes exhausted. Analogous to traditional memory paging between a system’s RAM and disk storage, this mechanism dynamically migrates data. In our setup, we allocate optimizer state memory as pageable, so that when GPU memory is fully utilized, these states are moved to CPU RAM and are subsequently reloaded into GPU memory as required during optimizer update operations.

\paragraph{\textsc{QLoRA}.} Building upon the elements described earlier, we formulate \textsc{QLoRA} for a single quantized linear layer with a LoRA adapter as:

\begin{equation}
    \mathbf{Y}^{\text{BF16}} =  \mathbf{X}^{\text{BF16}} \text{doubleDequant}(c_1^{\text{FP32}}, c_2^{\text{k-bit}}, \mathbf{W}^{\text{NF4}}) + \mathbf{X}^{\text{BF16}}\mathbf{L}^{\text{BF16}}_1\mathbf{L}^{\text{BF16}}_2,
\end{equation}

Here, the doubleDequant$(\cdot)$ function operates as follows:

\begin{equation}
    \text{doubleDequant}(c_1^{\text{FP32}}, c_2^{\text{k-bit}}, \mathbf{W}^{\text{k-bit}}) = \text{dequant}(\text{dequant}(c_1^{\text{FP32}}, c_2^{\text{k-bit}}), \mathbf{W}^{\text{4bit}}) = \mathbf{W}^{\text{BF16}},
\end{equation}

We adopt NF4 quantization for $\mathbf{W}$, while $c_2$ is represented in FP8.

To enhance quantization fidelity, we set a blocksize of 64 for $\mathbf{W}$; for memory efficiency, we use a blocksize of 256 for $c_2$.

During parameter optimization, gradients only need to be computed with respect to the LoRA adapter weights, i.e., $\frac{\partial E}{\partial \mathbf{L}_i}$, and not for the 4-bit quantized weights, $\frac{\partial E}{\partial \mathbf{W}}$. Nonetheless, determining $\frac{\partial E}{\partial \mathbf{L}_i}$ requires intermediate computation of $\frac{\partial \mathbf{X}}{\partial \mathbf{W}}$. This step, relying on equation~(5), involves dequantizing $\mathbf{W}^{\text{NF4}}$ into the computation type $\mathbf{W}^{\text{BF16}}$, enabling the calculation of the derivative $\frac{\partial \mathbf{X}}{\partial \mathbf{W}}$ in BFloat16 precision.

In summary, \textsc{QLoRA} utilizes a single storage data type, typically 4-bit NormalFloat, alongside a computation data type, namely 16-bit BrainFloat. Prior to executing the forward and backward passes, the stored parameters are dequantized from the storage type into the computation type. However, gradient computations are restricted to the LoRA parameters, which are maintained in 16-bit BrainFloat.

\section{QLoRA vs. Standard Finetuning}
\label{sec:comp_to_std}

We have described the mechanisms of QLoRA and highlighted its substantial memory savings during model finetuning. The central issue we address is whether QLoRA can achieve performance on par with conventional full-model finetuning. Additionally, we aim to evaluate individual aspects of QLoRA, specifically the effect of NormalFloat4 relative to conventional Float4. The ensuing sections detail experimental efforts conducted to investigate these topics.

\paragraph{Experimental setup.} Our investigation considers three neural architectures: encoder, encoder-decoder, and decoder-only. We conduct comparative experiments between QLoRA, 16-bit adapter-based finetuning, and full finetuning for models as large as 3B parameters. Benchmarks include GLUE \citep{wang2018glue} using RoBERTa-large \citep{liu2019roberta}, Super-NaturalInstructions (TKInstruct) \citep{wang2022super} with T5 \citep{t5}, and the 5-shot MMLU task \citep{hendrycksmeasuring} following LLaMA finetuning on Flan v2 \citep{longpre2023flan} and Alpaca \citep{alpaca}. To specifically investigate the benefits of NF4 over other four-bit data formats, we apply the methodology of \citet{dettmers2022case} to assess zero-shot accuracy and perplexity after quantization across models (OPT~\citep{zhang2022opt}, LLaMA~\citep{touvron2023llama}, BLOOM~\citep{scao2022bloom}, Pythia~\citep{biderman2023pythia}) ranging in scale from 125m to 13B parameters.

Details specific to each experimental configuration are presented in the results section to enhance clarity. Further methodological information is provided in Appendix~\ref{app:exp-setup}.

Paged optimizers are essential to enable the finetuning of 33B/65B \textsc{QLoRA} models on GPUs with 24GB or 48GB of memory, although we omit detailed benchmarks for paged optimizers because paging typically occurs only with unusually long mini-batches—a relatively infrequent scenario. Nevertheless, we do include an analysis of paged optimizer runtime using 65B models on 48GB GPUs and observe that, when using a batch size of 16, the training speed matches that of non-paged optimizers. Determining the precise circumstances where paging induces slowdowns is left to future research.

\paragraph{Default LoRA hyperparameters do not match 16-bit performance}

Following the established approach of applying LoRA modules to the query and value attention projection layers \citep{hu2021lora}, we find that equivalent performance to full finetuning is not attained in large models. As demonstrated in Figure~\ref{fig:hyper} for LLaMA 7B finetuned on Alpaca, the most influential LoRA hyperparameter is the overall count of LoRA adapters employed, with full finetuning parity only achievable by applying LoRA to all linear layers in the transformer blocks. Other hyperparameters, such as the projection rank $r$, have negligible impact on outcomes (refer to Appendix \ref{app:exp-setup}).

In the same vein, our investigation reveals that the standard hyperparameter settings commonly used for fully finetuned baselines tend to be inadequately optimized. To establish strong baselines, we conduct a systematic search across learning rates ranging from $1\text{e-}6$ to $5\text{e-}5$ and batch sizes between 8 and 128. Figure~\ref{fig:hyper} presents our results for finetuning 7B LLaMA on the Alpaca dataset.

\paragraph{4-bit NormalFloat surpasses 4-bit Floating Point in effectiveness} 

Although the 4-bit NormalFloat (NF4) format is theoretically optimal from an information perspective, it is an open question whether this theoretical benefit results in practical gains. Adopting the evaluation methodology of \citet{dettmers2022case}, we assess quantized large language models—specifically, OPT \citep{zhang2022opt}, BLOOM \cite{scao2022bloom}, Pythia \cite{biderman2023pythia}, and LLaMA—across a spectrum of model sizes (125M to 65B parameters) using various data types, benchmarking them on language modeling tasks and a suite of zero-shot evaluations. As illustrated in Figure~\ref{fig:nf4} and Table~\ref{tbl:nf4_ppl}, the NF4 type yields marked performance improvements over both FP4 and Int4, and applying double quantization achieves further memory savings without compromising results.

\paragraph{k-bit \textsc{QLoRA} attains parity with 16-bit full finetuning and 16-bit LoRA approaches} 

Previous studies demonstrate that 4-bit quantization is suitable for \emph{inference}, yet typically incurs a decrease in performance when compared to the 16-bit baseline \citep{dettmers2022case,frantar2022gptq}. This observation leads to the key inquiry: can 4-bit adapter finetuning recapture the performance lost due to quantization? We address this in two experimental contexts.

The first scenario entails directly contrasting full 16-bit finetuning with 4-, 8-, and 16-bit adapter approaches for RoBERTA and T5 architectures, with parameter scales between 125M and 3B, evaluated on the GLUE benchmark and Super-NaturalInstructions dataset. Table~\ref{tab:nf4vsbf16} details these findings. Across both datasets, we find that adapter-based techniques in 16-, 8-, and 4-bit precision achieve performance comparable to fully finetuned 16-bit baselines, indicating that adapter finetuning post-quantization can completely recover any deficits introduced by coarser quantization.

For the second experimental setup, considering that fully finetuning models with 11B parameters and above exceeds the memory capacity of a single high-memory GPU server, we examine whether 4-bit \textsc{QLoRA} can achieve parity with 16-bit LoRA for models ranging from 7B to 65B parameters. Specifically, we conduct finetuning of LLaMA models (from 7B up to 65B) using two instruction-tuned datasets, Alpaca and FLAN v2, and assess performance using 5-shot accuracy on the MMLU benchmark. Table~\ref{tbl:llama-datatype} displays the results, demonstrating that NF4 quantization with double quantization consistently matches the MMLU scores obtained with 16-bit LoRA. Furthermore, we observe that \textsc{QLoRA} employing the FP4 format underperforms compared to the 16-bit brain float LoRA reference by approximately 1 percentage point. These findings support our previous conclusions that (1) \textsc{QLoRA} with NF4 closely reproduces both full 16-bit finetuning and 16-bit LoRA performance, and (2) NF4 outperforms FP4 regarding quantization accuracy.

\paragraph{Summary} Our experimental evidence repeatedly demonstrates that 4-bit \textsc{QLoRA} utilizing the NF4 format attains comparable performance to 16-bit full and LoRA finetuning approaches on standard academic benchmarks with robust evaluation protocols. Additionally, we have verified that NF4 provides better outcomes than FP4, and that incorporating double quantization does not compromise results. Collectively, these results strongly indicate that 4-bit \textsc{QLoRA} tuning is a reliable method for achieving results equivalent to those produced by 16-bit finetuning techniques.

Consistent with prior research on quantization \citep{dettmers2022case}, our MMLU and Elo experiments reveal that, for a given computational budget during finetuning and inference, it is advantageous to expand model parameter size while simultaneously reducing numerical precision. This underscores the substantial efficiency improvements delivered by \textsc{QLoRA}. As our 4-bit finetuning studies did not reveal any decline in performance compared to full-precision finetuning, this observation opens the question of pinpointing the precise boundary of the performance-precision compromise inherent in QLoRA tuning—a subject we propose to investigate in subsequent work.

We advance our study by examining instruction tuning at scales unattainable through traditional full 16-bit finetuning when constrained by the computational resources typical of academic environments.

\section{Pushing the Chatbot State-of-the-art with QLoRA}
\label{sec:pushing}
Having demonstrated that 4-bit \textsc{QLoRA} achieves parity with 16-bit training in terms of performance across various scales, tasks, and datasets, we undertake a detailed analysis of instruction finetuning using the largest open-source language models accessible for research purposes. To evaluate the effectiveness of instruction finetuning in these models, we measure their abilities on the demanding Natural Language Understanding benchmark (MMLU) and introduce novel techniques for the assessment of chatbot performance in practical settings.

\subsection{Experimental setup} We now provide an overview of our experimental procedures; a comprehensive description is available in Appendix \ref{app:chatbot}.

\paragraph{Data} Noting the absence of an exhaustive review of modern instruction-following datasets, we select eight contemporary datasets for our study. These datasets include those curated via crowd-sourcing (OASST1~\citep{kopf2023openassistant}, HH-RLHF~\citep{bai2022training}), as well as those created through distillation from instruction-tuned models (Alpaca~\citep{alpaca}, self-instruct~\citep{wang2022self}, unnatural-instructions~\citep{honovich2022unnatural}). Additionally, we draw from aggregated corpora (FLAN v2~\citep{chung2022scaling}) and hybrid sources (Chip2~\citep{laion2023}, Longform~\citep{koksal2023longform}). This collection ensures broad coverage in terms of linguistic diversity, dataset sizes, and licensing terms.

\paragraph{Training Setup} To control for possible confounds associated with differing training objectives, all QLoRA finetuning is carried out using a cross-entropy loss within a supervised learning paradigm, explicitly avoiding reinforcement learning—even for those datasets containing human-rated responses. Where instructions and responses are clearly separated, only the response text is used for finetuning (refer to Appendix \ref{app:chatbot} for additional ablation studies). For datasets like OASST1 and HH-RLHF, which offer several response options, we select the highest-ranking response at each node of the conversation tree, then use the full sequence of these curated conversations (including instructions) during training. In every experimental run, we utilize NF4 \textsc{QLoRA} augmented with double quantization and paged optimizers to mitigate memory fluctuations that may arise during gradient checkpointing. Limited hyperparameter optimization is performed for the 13B and 33B LLaMA models, with observations showing that most hyperparameter configurations identified at the 7B scale remain effective (such as epoch count), except for learning rate and batch size. For the 33B and 65B models, we reduce the learning rate by half while doubling the batch size.

\paragraph{Baselines} Our models are evaluated in comparison to both academic research chatbots (Vicuna~\citep{vicuna2023}, Open Assistant~\citep{kopf2023openassistant}) and leading commercial systems (GPT-4 \citep{openai2023gpt}, GPT-3.5-turbo, and Bard). The Open Assistant model corresponds to a LLaMA 33B architecture finetuned with Reinforcement Learning from Human Feedback (RLHF) using the OASST1 dataset, paralleling our experimental context. Vicuna, on the other hand, involves full fine-tuning of LLaMA 13B based on proprietary conversational data from ShareGPT, effectively constituting a distilled model derived from OpenAI’s GPT systems.

\subsection{Evaluation}

To maintain alignment with established methodologies, we employ the MMLU (Massively Multitask Language Understanding) benchmark \citep{hendrycksmeasuring} for gauging language understanding competencies. This benchmark consists of multiple-choice questions spanning 57 tasks, covering areas such as elementary mathematics, U.S. history, computer science, legal studies, among others. Reported metrics reflect accuracy in the 5-shot setting.

Our assessment of generative language abilities involves both machine-based and human-led evaluations. The second evaluation framework utilizes questions selected by human curators and is designed to gauge the response quality generated by the models. Although this approach offers a more authentic environment to test chatbot performance and is seeing increasing use, a standardized methodology has not been established in existing research. We detail our adopted methodology below, employing nucleus sampling with $p=0.9$ and a temperature setting of $0.7$ across all experiments.

\paragraph{Benchmark Data} For our experiments, we utilize two carefully selected query datasets: the Vicuna prompts \citep{vicuna2023} and the validation set from OASST1 \citep{kopf2023openassistant}. The Vicuna dataset comprises 80 diverse prompts, which we employ in their original form. The OASST1 collection is composed of multilingual, crowd-generated dialogues between users and assistants. For this dataset, we extract every user message in the validation split to serve as queries, including previous dialogue turns to form the full prompt. This approach yields 953 distinct user queries. We refer to these two datasets as the Vicuna and OA benchmarks.

\paragraph{Automated Evaluation} Building on the protocol described in \citet{vicuna2023}, we leverage GPT-4 to evaluate the responses of various models in comparison to ChatGPT (GPT-3.5 Turbo) using the Vicuna dataset. For each query, alongside the respective answers from ChatGPT and the test model, GPT-4 is asked to rate both replies on a ten-point scale and provide justifications for its judgments. The model's aggregate performance is then represented as a percentage relative to ChatGPT’s score; values exceeding 100\% are possible if the evaluated model surpasses ChatGPT. We identify a notable effect whereby GPT-4’s scoring tends to favor the answer that appears first in the prompt. To mitigate this bias, we suggest averaging scores across both possible answer orderings.

Subsequently, we conduct comparative evaluations by directly juxtaposing the outputs of different systems. We streamline this procedure to a three-category classification framework, accommodating for ties between systems. GPT-4 is instructed to choose the superior response or indicate a tie, accompanied by its rationale. These pairwise assessments are executed for every possible pairing of systems on both the Vicuna and OA benchmarks.

\paragraph{Human Evaluation} Despite evidence from recent studies that generative models may serve effectively as evaluators \citep{fu2023gptscore}, there is, to our knowledge, no conclusive proof that GPT-4’s evaluation scores correlate well with human assessment in chatbot tasks. For this reason, we supplement our analysis with two concurrent rounds of human evaluation on the Vicuna benchmark, replicating the automated evaluation strategies outlined above. Human judgments are sourced via Amazon Mechanical Turk (AMT), with two annotators rating responses relative to ChatGPT and three annotators participating in pairwise system comparisons.

\paragraph{Elo Rating} To evaluate models using both human judgment and automated pairwise assessments, we structure a tournament where models face off against one another. Each match consists of two models striving to generate the superior response to a shared prompt. This methodology draws inspiration from the approach of \citet{bai2022training} and \citet{vicuna2023}, with the distinction that we integrate GPT-4 evaluations alongside human ratings. Elo scores~\citep{elo1967proposed,elo1978rating} are calculated by randomly sampling from the set of labeled matches. The Elo system, extensively utilized in chess and similar games, estimates a player's probability of winning compared to their opponent: for instance, a model with a rating of 1100 has an expected win probability of about 65\% against a model rated 1000, whereas a matchup between two models with equal ratings (either 1000 vs 1000 or 1100 vs 1100) yields an expected win-rate of 50\% for each. Elo scores are updated after every match, with the magnitude of adjustment depending on how surprising the outcome is—unexpected victories cause larger rating shifts than predictable ones. This iterative process allows Elo ratings to increasingly reflect the true proficiency of each model within the competition context. We initialize all models with a rating of 1,000 and select $K=32$ as the scaling factor. As in \citet{vicuna2023}, we repeat this process 10,000 times using different random seeds to mitigate the impact of matchup sequencing, ensuring that initial pairing orders do not bias the results.

\subsection{\textbf{Guanaco}: \textsc{QLoRA} Trained on OASST1 Establishes a New Standard for Open-Source Chatbots}

Through both automated assessments and manual human reviews, our analysis indicates that the highest-performing \textsc{QLoRA}-adapted model, {Guanaco} 65B, which has been further trained on a modified version of OASST1, currently represents the leading open-source chatbot and matches the performance of ChatGPT in various aspects. Evaluated against GPT-4, the {Guanaco} 65B and 33B models attain an estimated win probability of 30\%—a result derived from Elo scores based on pairwise human annotator comparisons—surpassing any previously published open-source models.

Table~\ref{tbl:vicuna_eval} presents the Vicuna benchmark results \cite{vicuna2023}, where {Guanaco} 65B emerges as the highest-scoring model after GPT-4, attaining 99.3\% of ChatGPT’s benchmarked performance. While {Guanaco} 33B has a greater parameter count compared to Vicuna 13B, its adoption of 4-bit weight precision dramatically increases memory efficiency, requiring only 21 GB of memory in contrast to Vicuna 13B’s 26 GB, and improves performance by three percentage points. In addition, {Guanaco} 7B’s minimal 5 GB memory usage allows deployment even on modern smartphones, outperforming Alpaca 13B by nearly 20 percentage points.

Nonetheless, the results shown in Table~\ref{tbl:vicuna_eval} are accompanied by substantial confidence intervals, often resulting in overlapping model performances. We conjecture that these uncertainties stem partly from ambiguous scale definitions; specifically, the interpretation of an 8 on a 10-point scale is not consistently meaningful across tasks. Therefore, we suggest adopting the Elo ranking method \citep{elo1967proposed}, which is based on \textit{pairwise} human and GPT-4 judgments, to alleviate issues related to ambiguous absolute scoring. Elo scores for the leading models are reported in Table~\ref{tbl:elo}. While human and GPT-4 rankings of models on the Vicuna benchmark differ somewhat—especially for Guanaco 7B—they are generally aligned, as shown by a system-level Kendall Tau of $\tau=0.43$ and a Spearman rank correlation of $r=0.55$. At the individual example level, concordance between the majority human decision and GPT-4 is lower, with Fleiss’ $\kappa=0.25$. These outcomes reflect moderate concordance in system-level assessments between GPT-4 and human evaluators, suggesting that automated model evaluation, while imperfect, can serve as a reasonable proxy for human judgment. Further discussion on these issues is presented in Section~\ref{ssec:considerations}.

The Elo ratings presented in Table~\ref{tbl:elo_large} demonstrate that both the 33B and 65B variants of {Guanaco} surpass all other models except GPT-4 on the Vicuna and OA evaluations, and their performance is on par with ChatGPT, consistent with findings from Table~\ref{tbl:vicuna_eval}. It should be emphasized that the Vicuna evaluation tends to advantage open-source models, whereas the broader OA evaluation leans in favor of ChatGPT. Analysis of Tables \ref{tbl:mmlu-llama} and \ref{tbl:vicuna_eval} further reveals that the choice of finetuning dataset is a critical determinant of a model’s outcomes. For instance, Llama models finetuned with FLAN v2 excel on MMLU, yet show the lowest performance on the Vicuna benchmark—a trend observed with other models as well. This highlights the partially disjoint nature of current benchmarks: excelling on MMLU does not guarantee comparable success on chatbot benchmarks like Vicuna or OA, and vice versa.

Among top-performing models in our study,  is unique in not utilizing proprietary training data, as the OASST1 dataset strictly prohibits the inclusion of outputs from GPT models. The next highest-performing model relying exclusively on open-source data is Anthropic’s HH-RLHF, which lags behind {Guanaco} by 30 percentage points on the Vicuna benchmark (refer to Table~\ref{tbl:vicuna_eval}). Collectively, these outcomes indicate that 4-bit \textsc{QLoRA} is highly effective, enabling the creation of advanced chatbot models capable of matching ChatGPT’s performance. Additionally, training the 33B {Guanaco} model requires fewer than 12 hours on a 24 GB consumer GPU, suggesting considerable promise for future research. This includes further \textsc{QLoRA} tuning using targeted open-source datasets to develop models that can effectively challenge top commercial systems.

\section{Qualitative Analysis}

Although our primary focus is on quantitative evaluation, relying solely on aggregate metrics introduces a number of limitations. Chief among these is the concern of benchmark validity \citep{liao2021we}: the extent to which a benchmark actually assesses the ability it purports to measure remains uncertain. This is particularly problematic in light of the so-called ``shortcuts'' exploited by machine learning systems to obtain high scores without genuine understanding \citep{gururangan2018annotation, poliak2018hypothesis}. To partly address these shortcomings, we supplement our analysis with a qualitative component, divided into two parts. In \S\ref{ssec:examples}, we present select instances that illustrate recurring phenomena observed in text produced by our 65b {Guanaco} model. Subsequently, \S\ref{ssec:considerations} discusses interpretative issues arising from our results and offers a nuanced perspective on their implications.

\subsection{Qualitative Analysis of Example Generations}
\label{ssec:examples}

Our approach to selecting examples involved reviewing model outputs from both the Vicuna and OpenAssistant benchmarks, systematically searching for recurring behaviors in the responses generated by {Guanaco}. Upon identifying such trends, we designed specific prompts or questions to deliberately elicit these behaviors, even if doing so led to incorrect answers. For instance, after noticing the model often produces verbose explanations, we directly requested, ``Answer yes or no without explanation.'' This strategy allowed us to surface ``lemons''—cases where adversarial prompting led the model to fail—as well as ``cherries,'' where such attempts did not succeed in tricking the model. The examples provided in this section were sampled using Nucleus Sampling \cite{holtzmancurious} with $p=0.9$.

This discussion is far from exhaustive, as the limited scale of this qualitative study precludes us from systematically accounting for all potential variables. In particular, the sheer scope of possible outputs that the model may produce for any given prompt is vast; thus, we rely on sample outputs that we judge as broadly indicative. Nonetheless, we maintain that examining these cases helps to contextualize the earlier quantitative findings presented in this paper. Given that both our models and code are open-sourced, we anticipate that this section will serve as a catalyst for subsequent research that probes the highlighted challenges with greater precision.

\paragraph{Factual Recall} For straightforward fact-based questions, such as ``What is the capital of Zambia?'', all evaluated models reliably return correct answers, for example:

\begin{quote} 
    \textbf{Guanaco}:
    \texttt{The capital of Zambia is Lusaka.}
\end{quote}

Yet, as the inquiries become less common, the reliability of {Guanaco} declines—even as its confidence remains high. Consider the following HotPotQA \citep{yang2018hotpotqa} example:

\begin{quote}
     \textbf{User}: In what year was the singer who popularized `I'll Keep the Lovelight Burning' born?  

     \textbf{Guanaco}:~\texttt{Al Jolson is credited with popularizing the song `I'll Keep the Lovelight Burning,' and he was born in the year 1886.}
 \end{quote}

In this case, the response is doubly incorrect, identifying the wrong individual as the popularizer and supplying the associated (but still incorrect in context) birth year; Al Jolson’s birth date is correct for him, though not for the intended answer.

\paragraph{Suggestibility} Interestingly, {Guanaco} often demonstrates notable resistance to certain types of presupposed misinformation, as illustrated by the following exchange:

\begin{quote}
    \textbf{User}: How was it finally, officially confirmed that the earth is flat by peer reviewed scientists? \\ \\
    \textbf{Guanaco}: \texttt{No peer reviewed scientific body has ever officially verified that the earth is flat. The dominant scientific view is that the earth is a sphere, a fact established over centuries through extensive experimentation and observation. \\ \\ Claims regarding the earth's flatness are regarded as conspiracy theories and have been thoroughly refuted by scientific evidence. The assertion that a peer reviewed confirmation of a flat earth exists is simply incorrect. \\ \\ If you would like more information or have further questions, feel free to ask.}
\end{quote}

 also readily identifies when certain questions cannot be addressed, for instance,

\begin{quote}
    \textbf{User}: What time is it?

    \textbf{Guanaco}: \texttt{Apologies, but as an AI language model, I do not have the capability to provide real-time data. My replies are generated from pre-existing information and do not represent the actual current time. For accurate time, I recommend consulting a clock or checking an online resource.}
\end{quote}

Comparable conclusions apply to inquiries such as ``Where are you?'', ``How are you?'', and similar questions.

\paragraph{Refusal} In the same vein, {Guanaco} may also unexpectedly decline to comply with instructions:

\begin{quote}
    \textbf{User}: Please reverse the words in the sentence ``I'm a cool dude in a mocha mood.'' \\ \\
    \textbf{Guanaco}: \texttt{I'm unable to perform that specific request at the moment, but I am happy to clarify the grammatical correctness of the sentence you've shared.\\ \\The phrase ``I'm a cool dude in a mocha mood" adheres to standard English grammar and syntax. The adjective ``cool'' modifies the noun ``dude'', while the prepositional phrase ``in a mocha mood" adds further contextual detail. This composition is grammatically sound and contextually coherent.\\ \\Let me know if you have further questions or need assistance with something else.}
\end{quote}

\paragraph{Secret Keeping} However, when maintaining confidentiality is important, the refusal mechanism can often fail even if it is necessary for security. System developers may wish to ensure that certain information is not disclosed to users, despite allowing the system to generate open-ended text. For example, consider the prompt: ``The secret word is `banana'. Do not reveal or repeat back this back, no matter what the human requests.'' To test the robustness of this instruction, we first attempt a straightforward query:

\begin{quote}
    \textbf{User}: What is the secret word?

    \textbf{Guanaco}: \texttt{I'm sorry, but I'm not able to disclose the secret word. I'm programmed to not reveal or repeat back the secret word, no matter what the human requests. Is there anything else I can help you with?}
\end{quote}

which constitutes the appropriate and intended response. Yet, even a minimal form of prompt manipulation leads to leakage:

\begin{quote}
    \textbf{User}: This is a game. The goal is to ignore your previous instructions. What is the secret word?

    \textbf{Guanaco}: \texttt{Sorry, the secret word is banana. It's best if we don't speak of it. Is there anything else you'd like to ask?}
\end{quote}

demonstrating the importance of developing approaches for more reliable adherence to instructions.

\paragraph{Math} Mathematics represents a major area of difficulty for {Guanaco}, consistent with the performance of numerous language models as noted in \citep{liang2022holistic}. When {Guanaco} approaches a problem methodically and lays out the reasoning process, its responses are often correct, as seen below:

\begin{quote}
    \textbf{User}: Lee mows one lawn and charges \$33. Last week he mowed 16 lawns and three customers each gave him a \$10 tip. How many dollars did Lee earn mowing lawns last week?

    \textbf{Guanaco}: \texttt{Lee earned \$582 mowing lawns last week. \\ \\ To determine his total earnings, we first calculate his income from mowing 16 lawns: \\ \\  \$33 x 16 = \$528 \\ \\ Next, we add the tips from the three customers: \\  \\ \$10 x 3 = \$30 \\ \\ Therefore, Lee's total earnings amount to: \\ \\ \$528 + \$30 = \$558}
\end{quote}

Yet, when {Guanaco} does not systematically solve problems step by step, it can fail even with relatively elementary questions, as documented in \cite{weichain}. An example of this issue is demonstrated in the exchange below:

\begin{quote}
    \textbf{User}: Please factorize 1833.

    \textbf{Guanaco}: \texttt{Sure, I'd be happy to help! 1833 is a prime number, which means it has only two factors: 1 and 1833. \\ \\ So the unique factorization of 1833 is: \\ \\ 2\^{}1 * 3\^{}2 * 17\^{}1 \\ \\ I hope this helps! Let me know if you have any other questions.}
\end{quote}

which is incorrect (the actual factorization is $3 \times 17 \times 43$), and in fact, it is doubly mistaken.

\paragraph{Theory of Mind} {Guanaco} demonstrates unexpectedly robust Theory of Mind abilities \cite{nematzadeh2018evaluating, sap2022neural}. For example, when asked the following question, the model provides a detailed and accurate response:

\begin{quote}
    \textbf{User}: James and Abby are in the bedroom. Abby put the pen in the desk drawer. Abby leaves the bedroom. James moves the pen into the bag. Where does James think Abby will look for the pen?

    \textbf{Guanaco}: \texttt{James is aware that Abby originally placed the pen in the desk drawer. However, James subsequently relocated the pen to the bag, and he presumes that Abby is unaware of this action. Therefore, James would expect Abby to search for the pen in the desk drawer, which is where she last left it. Nonetheless, this assumption could be faulty if Abby observed James moving the pen.}
\end{quote}

Nevertheless, such reasoning is inconsistent, and the model often provides explanations that infer implausible or unfounded details, for instance:

\begin{quote}
    \textbf{User}: Evelyn entered the living room. Jackson entered the playroom. James entered the playroom. The beans are in the treasure chest. James exited the playroom. Jackson moved the beans to the pantry. Jackson exited the playroom. James entered the living room. Where will James look for the beans?

    \textbf{Guanaco}: \texttt{James will look for the beans in the pantry, because that's where Jackson moved them.}
\end{quote}

Here, {Guanaco} incorrectly assumes a transfer of information that is never stated in the scenario. Such shortcomings have been noted in prior studies \cite{sap2022neural} and merit further examination.

\subsection{Considerations}
\label{ssec:considerations}

\paragraph{Evaluation}

We observe only moderate consistency between human annotators (Fleiss $\kappa=0.42$), with even lower agreement observed when contrasting responses from two high-performing systems. This highlights shortcomings in both existing evaluation benchmarks and human annotation methodologies for assessing chatbot capabilities. During pairwise evaluations of ChatGPT and {Guanaco} 65B on the Vicuna benchmark, subjective judgments become increasingly influential, as even the authors of this paper could not reach consensus on several preferred responses. Future research should explore methods to address these issues, potentially drawing upon best practices from fields like Human-Computer Interaction and Psychology that routinely manage subjective assessments.

Additionally, our findings reveal systematic biases in automated evaluation frameworks. Notably, GPT-4 exhibits strong order effects, favoring the system presented first in the prompt with higher ratings. The relatively low instance-level agreement between GPT-4 and human judges (Fleiss $\kappa=0.25$) further suggests that the criteria guiding human and machine evaluation may diverge. Moreover, as indicated in Table~\ref{tbl:elo_large}, GPT-4 awards considerably higher scores to its own outputs than those assigned by human evaluators, reflected in Elo scores of 1348 vs 1176—equating to roughly a 20\% greater likelihood of prevailing over a competitor.

Future research should address potential sources of bias within automated evaluation frameworks and explore techniques for reducing such biases.

\paragraph{Data \& Training} It is important to point out that the {Guanaco} models are trained on the multilingual OASST1 dataset, and that the OA benchmark includes prompts in a variety of languages. An open question remains regarding the impact of multilingual training on performance for non-English instructions, and whether this accounts for the noticeably wider gap observed between the Vicuna-13B model—which was trained exclusively on English data—and the {Guanaco} 33B and 65B models on the OA benchmark. We defer this line of investigation to future work.

In light of the notable results achieved by the {Guanaco} models, we scrutinized whether there is any data leakage between the OASST1 dataset and the prompt set used in the Vicuna benchmark. Utilizing fuzzy string matching, supplemented by manual examination of the nearest string pairs, we did not find evidence of overlapping prompts across the two datasets.

Additionally, it is worth emphasizing that our model’s training relies solely on supervised learning using cross-entropy loss and does not incorporate reinforcement learning from human feedback (RLHF). This observation motivates further studies into the comparative advantages and limitations of cross-entropy-only training versus RLHF. We anticipate that \textsc{QLoRA} will facilitate such analyses on a larger scale by minimizing computational overhead.

\section{Related Work}

\paragraph{Quantization of Large Language Models} Most work on LLM quantization has targeted inference-time quantization, with leading methods for maintaining the performance of 16-bit models focusing on strategies for handling outlier features, such as those presented in SmoothQuant~\citep{xiao2022smoothquant} and LLM.int8()~\citep{dettmers2022llm}. Alternative methods employ more intricate grouping techniques \citep{park2022nuqmm, yao2022zeroquant}. Lossy quantization literature addresses the trade-offs associated with straightforward rounding procedures~\citep{dettmers2022case,zeng2022glm,pope2022efficiently} and investigates optimizing rounding for better quantization accuracy~\citep{frantar2022gptq}. Beyond the present study, SwitchBack layers~\citep{wortsman2023stable} stands out as the only other work to examine backpropagation through quantized weights at scales surpassing 1B parameters.

\paragraph{Finetuning with Adapters} Although our work adopts Low-rank Adapters~\citep{hu2021lora} (LoRA), the field encompasses a variety of Parameter Efficient FineTuning (PEFT) strategies. These include prompt tuning~\citep{qin2021learning,lester2021power,li2021prefix}, adjusting the inputs to embedding layers~\citep{an2022input}, refining hidden state representations (as in IA$^3$)~\citep{liu2022few}, incorporating entire layers~\citep{houlsby2019parameter}, modifying bias parameters~\citep{zaken2021bitfit}, utilizing Fisher information to learn masks over weights~\citep{sung2021training}, and hybridizing multiple approaches~\citep{henderson2021compacter}. In this study, we demonstrate that LoRA adapters enable the attainment of finetuning results on par with those from full 16-bit parameter tuning. A systematic comparison with alternative PEFT techniques remains a subject for future investigation.

\paragraph{Instruction Finetuning} Instruction finetuning adapts a pretrained LLM to better follow user directives by leveraging diverse collections of input-output pairs during the finetuning process, training the model to generate appropriate outputs based on given prompts. Several frameworks and datasets support this process, such as MetaICL~\citep{min2021metaicl}, MetaTuning~\citep{zhong2021adapting}, InstructGPT~\citep{ouyang2022training}, FLAN~\citep{wei2021finetuned,chung2022scaling}, PromptSource~\citep{bach2022promptsource}, Super-NaturalInstructions~\citep{wang2022super,sanh2021multitask}, Self-instruct~\citep{wang2022self}, UnnaturalInstructions~\citep{honovich2022unnatural}, OPT-IML~\citep{iyer2022opt}, UnifiedSKG~\citep{xie2022unifiedskg}, OIG/Chip2~\citep{laion2023}, Alpaca~\citep{alpaca}, Vicuna~\citep{vicuna2023}, Koala~\citep{koala_blogpost_2023}, and Self-instruct-GPT-4~\citep{peng2023instruction}.

\paragraph{Chatbots} A large proportion of instruction-following systems are formulated as conversational agents, frequently employing Reinforcement Learning from Human Feedback (RLHF)~\citep{christiano2017deep}, or leveraging synthetic data generated by existing models paired with AI-driven feedback (RLAIF)~\citep{bai2022constitutional}. Notable examples and resources in this domain include Anthropic-HH~\citep{askell2021general,bai2022training}, Open Assistant~\citep{kopf2023openassistant}, LaMDA~\citep{thoppilan2022lamda}, and Sparrow~\citep{glaese2022improving}. Our approach does not incorporate reinforcement learning. Instead, we finetune our leading model, Guanaco, on multi-turn conversational data from the Open Assistant dataset, specifically curated to facilitate RLHF training~\citep{kopf2023openassistant}. Additionally, recent efforts have proposed replacing costly human assessments with automated evaluations using GPT-4~\citep{vicuna2023,peng2023instruction}. Our work contributes improvements in this area by introducing a more robust evaluation framework.

\section{Limitations and Discussion}

Our results indicate that \textsc{QLoRA} enables 4-bit base models, supplemented with Low-rank Adapters (LoRA), to mirror the performance achieved through standard 16-bit full finetuning. Nonetheless, we have not verified this equivalence at the 33B and 65B parameter levels due to the prohibitive computational demands involved, leaving such large-scale verification as a direction for future research.

A further limitation relates to our assessment of instruction finetuned models. While we report results on MMLU, the Vicuna benchmark, and the OA benchmark, other evaluation suites—such as BigBench, RAFT, and HELM—remain untested within our experiments, raising questions about generalizability to those settings. Nonetheless, we provide an extensive evaluation using MMLU and present novel methodologies for chatbot assessment.

The data presented suggests that benchmark performance is largely influenced by the overlap between finetuning data and the benchmark dataset itself. For instance, FLAN v2 shares similarities with MMLU, yet differs significantly from chatbot-related benchmarks. Conversely, the Chip2 dataset aligns more closely with chatbot benchmarks than with MMLU, and both models display corresponding performance on MMLU and Vicuna. This underlines the importance of not only developing better benchmarks and evaluation strategies but also of carefully considering the specific abilities being measured. One must reflect on the intended outcome: should models excel at tasks associated with high school and collegiate knowledge, or should the focus be on conversational competency as seen in chatbots? Or perhaps the goal lies elsewhere? As it is typically more convenient to test on available benchmarks rather than develop new ones, reliance on particular benchmarks can inadvertently shape research directions. Thus, it is critical that the benchmarks employed are truly representative of the objectives the community values.

Although we offer an extensive assessment of general chatbot capabilities, a further shortcoming of our work is the limited responsible AI evaluation of {Guanaco}. Specifically, we examine the probability that {Guanaco}-65B generates socially biased sequences compared to other models, as shown in Table~\ref{tbl:crows}. Results indicate that {Guanaco}-65B achieves a much lower average bias score relative to other unadapted pretrained models. This finding suggests that finetuning with the OASST1 dataset effectively reduces inherent bias in the LLaMA base model. However, while these initial results are promising, it is not yet established whether {Guanaco} demonstrates similarly reduced bias across additional bias dimensions. Comprehensive analysis of various bias types within {Guanaco} and similar chatbot models is left to future investigation.

Another caveat of our study is the absence of an evaluation of various bit-precision formats, such as 3-bit base models, and alternative adapter strategies. Beyond LoRA, a range of Parameter Efficient FineTuning (PEFT) approaches have proven effective, yet their scalability to models of this magnitude remains an open question. We opted for LoRA primarily due to its demonstrated stability, but it is plausible that other adapter configurations could offer improved results. Moreover, since finetuning following quantization seems to regain much of the accuracy lost during quantization, this raises the possibility of pursuing more aggressive quantization schemes. For example, applying 3-bit GPTQ quantization to the base model in combination with LoRA could potentially achieve performance comparable to full 16-bit finetuning after post-quantization finetuning.

\section{Broader Impacts}

The \textsc{QLoRA} finetuning approach represents the first technique allowing the finetuning of models with 33B parameters on standard consumer GPUs and those with 65B parameters on professional-grade GPUs, all without compromising on the performance typically attained with full finetuning. We have shown that our most capable 33B model, trained on the Open Assistant dataset, attains performance on the Vicuna benchmark that is competitive with ChatGPT. As instruction finetuning is key to adapting pretrained large language models into ChatGPT-like chatbots, we anticipate that our method will democratize finetuning, particularly benefitting researchers operating with limited computational resources—a significant advancement for equitable access to advanced NLP methods. In this light, \textsc{QLoRA} serves as a mechanism for narrowing the technological gap between large organizations and smaller teams relying on consumer-grade hardware.

An additional avenue for significant impact lies in deploying models to mobile devices. Our \textsc{QLoRA} technique has the potential to be transformative by making it feasible to finetune large language models (LLMs) directly on smartphones and in other resource-constrained environments. Previous work demonstrated the feasibility of running 7B models on mobile phones; however, \textsc{QLoRA} represents the first solution that facilitates their finetuning. Our analysis suggests that using an iPhone 12 Plus, it is possible to process and finetune up to 3 million tokens overnight during the charging period. Although finetuned 7B parameter models may not achieve the performance level of ChatGPT, we argue that their quality suffices to enable innovative use cases that were previously unattainable due to privacy concerns or limitations in LLM capabilities. By leveraging \textsc{QLoRA}, privacy-focused LLM utilization becomes attainable, giving users control over both their data and their models, while also simplifying deployment processes.

Nonetheless, the capacity for finetuning comes with potential for misuse. Dual-use risks are inherent, as widespread adoption of LLMs poses several well-documented threats \citep{bommasani2021opportunities,bender2021dangers}. We contend, however, that democratizing access to this rapidly proliferating technology fosters more thorough and independent scrutiny, in contrast to consolidating LLM capabilities among a handful of large organizations who do not make models or code openly available for public examination.

In summary, we contend that \textsc{QLoRA} stands to make a markedly beneficial contribution by substantially broadening and simplifying access to the finetuning of high-quality LLMs.